% ==============================================================================
% CONFIG: PACKAGES.TEX
% Các thư viện LaTeX sử dụng cho cuốn sách "Học Lập Trình Web - Chuyên Sâu Frontend"
% Trình biên dịch: XeLaTeX (Hỗ trợ tiếng Việt Unicode & Font hệ thống)
% ==============================================================================

% 0. Bộ lọc triệt tiêu toàn bộ cảnh báo nhật ký (Silence Harmless Warnings)
\usepackage{silence}
\WarningFilter{tcolorbox}{Using nobreak failed}
\WarningFilter{babel}{No hyphenation patterns}
\WarningFilter{microtype}{Unable to apply patch}
\WarningFilter{rerunfilecheck}{File `main.out'}

% 1. Hỗ trợ phông chữ Unicode & Tiếng Việt cho XeLaTeX
\usepackage{fontspec}
\usepackage[vietnamese]{babel}
\shorthandoff{"}

% 2. Định dạng lề trang sách chuyên nghiệp (Chuẩn khổ A4 / Book)
\usepackage[
    paper=a4paper,
    top=2.6cm,
    bottom=2.6cm,
    left=2.5cm,   % Lề trái cân đối
    right=2.5cm,  % Lề phải cân đối
    headheight=26pt,
    headsep=20pt,
    footskip=30pt
]{geometry}

% 3. Bảng màu cao cấp (Modern Palettes)
\usepackage{xcolor}
\usepackage{colortbl}
\definecolor{primaryBlue}{HTML}{0F4C81}     % Classic Navy/Royal Blue (Màu chủ đạo)
\definecolor{secondaryTeal}{HTML}{008080}   % Teal Accent
\definecolor{darkText}{HTML}{1A1A1A}        % Màu chữ đen dịu
\definecolor{accentOrange}{HTML}{E65100}    % Cam điểm nhấn
\definecolor{accentPurple}{HTML}{5E35B1}    % Tím sang trọng
\definecolor{codeBg}{HTML}{F8FAFC}          % Nền code xám nhạt chuẩn VS Code Light
\definecolor{codeBorder}{HTML}{CBD5E1}      % Viền code Slate
\definecolor{codeComment}{HTML}{64748B}     % Màu chú thích code Slate
\definecolor{codeKeyword}{HTML}{D97706}     % Từ khóa Cam/Hổ phách đậm
\definecolor{codeString}{HTML}{059669}      % Chuỗi ký tự Xanh Ngọc Emerald
\definecolor{codeKey}{HTML}{0284C7}         % Tên thuộc tính/Headers Xanh Cyan
\definecolor{codeNumber}{HTML}{7C3AED}      % Giá trị số Tím Indigo

% Màu cho các loại Callout Boxes
\definecolor{noteBg}{HTML}{F0F7FF}
\definecolor{noteBorder}{HTML}{0F4C81}
\definecolor{tipBg}{HTML}{F0FDF4}
\definecolor{tipBorder}{HTML}{16A34A}
\definecolor{warnBg}{HTML}{FFFBEB}
\definecolor{warnBorder}{HTML}{D97706}
\definecolor{conceptBg}{HTML}{F5F3FF}
\definecolor{conceptBorder}{HTML}{7C3AED}

% 4. Đồ họa, Vẽ hình TikZ & Chèn ảnh
\usepackage{graphicx}
\usepackage{float}
\usepackage{caption}
\usepackage{subcaption}
\usepackage{tikz}
\usetikzlibrary{shapes, shadows, arrows.meta, positioning, calc, decorations.pathreplacing, fit}

% 5. Bảng biểu nâng cao
\usepackage{tabularx}
\usepackage{booktabs}
\usepackage{multirow}
\usepackage{array}
\usepackage{ragged2e} % Hỗ trợ căn lề RaggedRight thông minh triệt tiêu lỗi giãn chữ thưa trong tcolorbox

% 6. Hộp ghi chú & Callout cao cấp (tcolorbox)
\usepackage[most]{tcolorbox}

% 7. Tô màu mã nguồn (Listings)
\usepackage{listings}

% 8. Header & Footer chuyên nghiệp
\usepackage{fancyhdr}

% 9. Định dạng tiêu đề Chapter, Section & Part (titlesec & titletoc)
\usepackage{titlesec}
\usepackage{titletoc}

% 10. Hyperlink & Metadata PDF
\usepackage{hyperref}
\hypersetup{
    colorlinks=true,
    linkcolor=primaryBlue,
    citecolor=secondaryTeal,
    urlcolor=accentOrange,
    pdfauthor={Đội Ngũ Lập Trình Web},
    pdftitle={Học Lập Trình Web - Chuyên Sâu Frontend},
    pdfkeywords={Frontend, HTML5, CSS3, JavaScript, Web Development, XeLaTeX}
}

% 11. Tiện ích danh sách, Căn lề & Giãn dòng chuẩn sách in
\usepackage{indentfirst} % Bắt buộc lùi đầu dòng cho mọi đoạn văn (kể cả đoạn đầu sau tiêu đề)
\setlength{\parindent}{18pt} % Độ lùi đầu dòng 18pt
\setlength{\parskip}{2pt}    % Khoảng cách giữa các đoạn văn vừa phải

\usepackage{enumitem}
\usepackage[normalem]{ulem}
\setlist[itemize]{leftmargin=1.8em, itemsep=2pt, parsep=0pt, topsep=3pt, partopsep=0pt}
\setlist[enumerate]{leftmargin=1.8em, itemsep=2pt, parsep=0pt, topsep=3pt, partopsep=0pt}
% \usepackage{microtype}
% \microtypesetup{patch=none}
\usepackage{setspace}
\setstretch{1.15} % Giãn dòng 1.15 gọn gàng, tiết kiệm trang sách
\usepackage{needspace}
